% year: תשפ"ה
% type: בוחן
% semester: א
% source: exams/בחנים/תשפ״ה, בוחן, 9141.pdf
% part: א
% number: 1
% points: 20
% categories: func-limits

\begin{question}
חשבו את הגבול הבא:
\[
\lim_{x\to\infty}\left(\sqrt{x^2+5x}-\sqrt{x^2-5x}\right)
\]
\end{question}

\begin{hint}
זהו גבול מהצורה $\infty-\infty$. כפלו וחלקו בצמוד $\sqrt{x^2+5x}+\sqrt{x^2-5x}$.
\end{hint}

\begin{solution}
עבור $x>5$ שני השורשים מוגדרים וחיוביים. זהו גבול מהצורה $\infty-\infty$, ולכן נכפול ונחלק בביטוי הצמוד (החיובי):
\[
\sqrt{x^2+5x}-\sqrt{x^2-5x}
=\frac{(x^2+5x)-(x^2-5x)}{\sqrt{x^2+5x}+\sqrt{x^2-5x}}
=\frac{10x}{\sqrt{x^2+5x}+\sqrt{x^2-5x}} .
\]
נחלק מונה ומכנה ב-$x$. מכיוון ש-$x>0$ מתקיים $x=\sqrt{x^2}$, ולכן
\[
=\frac{10}{\sqrt{1+\frac5x}+\sqrt{1-\frac5x}} .
\]
כאשר $x\to\infty$: $\frac5x\to 0$, ומרציפות פונקציית השורש ב-1 נקבל $\sqrt{1\pm\frac5x}\to 1$. לפי אריתמטיקה של גבולות (המכנה שואף ל-$2\neq 0$):
\[
\lim_{x\to\infty}\left(\sqrt{x^2+5x}-\sqrt{x^2-5x}\right)=\frac{10}{1+1}=\boxed{5}.
\]
\end{solution}
